\subsection{Example: Arithmetic Operations for Rational Numbers}
\label{Section 2.1.1}

Suppose we want to do arithmetic with rational numbers.  We want to be able to
add, subtract, multiply, and divide them and to test whether two rational
numbers are equal.

Let us begin by assuming that we already have a way of constructing a rational
number from a numerator and a denominator.  We also assume that, given a
rational number, we have a way of extracting (or selecting) its numerator and
its denominator.  Let us further assume that the constructor and selectors are
available as procedures:

\begin{itemize}
\item
\code{make\_rat(}\meta{n}\code{, }\meta{d}\code{)}
returns the rational number whose numerator is the integer
\meta{n} and whose denominator is the integer \meta{d}.

\item
\code{numer(}\meta{x}\code{)} returns the numerator of the rational number
\meta{x}.

\item
\code{denom(}\meta{x}\code{)} returns the denominator of the rational number
\meta{x}.
\end{itemize}

\noindent
We are using here a powerful strategy of synthesis: \newterm{wishful thinking}.
We haven't yet said how a rational number is represented, or how the procedures
\code{numer}, \code{denom}, and \code{make\_rat} should be implemented.  Even
so, if we did have these three procedures, we could then add, subtract,
multiply, divide, and test equality by using the following relations:

\begin{align*}
  \frac{n_1}{d_1} + \frac{n_2}{d_2} 	&= \frac{n_1 d_2 + n_2 d_1}{d_1 d_2}, \\
  \frac{n_1}{d_1} - \frac{n_2}{d_2} 	&= \frac{n_1 d_2 - n_2 d_1}{d_1 d_2}, \\
  \frac{n_1}{d_1} \cdot \frac{n_2}{d_2}	&= \frac{n_1 n_2}{d_1 d_2}, \\
  \frac{{n_1 / d_1}}{{n_2 / d_2}} 	&= \frac{n_1 d_2}{d_1 n_2}, \\
  \frac{n_1}{d_1} 			&= \frac{n_2}{d_2} \quad
						{\rm\ if\ and\ only\ if\quad}
						n_1 d_2 = n_2 d_1.
\end{align*}
\noindent
We can express these rules as procedures:

\begin{codeBlock}
>>> def add_rat(x, y):
...     return make_rat(numer(x) * denom(y) + numer(y) * denom(x),
...                     denom(x) * denom(y))

>>> def sub_rat(x, y):
...     return make_rat(numer(x) * denom(y) - numer(y) * denom(x),
...                     denom(x) * denom(y))

>>> def mul_rat(x, y):
...     return make_rat(numer(x) * numer(y),
...                     denom(x) * denom(y))

>>> def div_rat(x, y):
...     return make_rat(numer(x) * denom(y),
...                     denom(x) * numer(y))

>>> def is_equal_rat(x, y):
...     return numer(x) * denom(y) == numer(y) * denom(x)
\end{codeBlock}

Now we have the operations on rational numbers defined in terms of the selector
and constructor procedures \code{numer}, \code{denom}, and \code{make\_rat}.
But we haven't yet defined these.  What we need is some way to glue together a
numerator and a denominator to form a rational number.

\subsubsection*{Pairs}

To enable us to implement the concrete level of our data abstraction, we need a
way of holding two things together as one.  Python's \newterm{tuple} does this:
several values kept in order, written one after another with commas between
them.  It is the comma that makes the tuple; the parentheses usually put round
it are for grouping, as they are in arithmetic, and may be left out where
nothing is ambiguous.\footnote{Which is a place where Python will let you write
something you did not mean.  Since it is the comma and not the brackets that
does the work, \code{(1)} is merely the number 1 with brackets round it, and a
tuple of one must be written with a comma you would not otherwise think to
put---\code{(1,)}.  Worse, a comma left at the end of a line is perfectly legal:
\code{x = 5,} binds \code{x} to the one-tuple \code{(5,)} and gives no
sign of having done so, and the mistake is generally discovered some distance
away, when something expecting a number is handed a tuple.  The empty tuple
\code{()} is the one case where the brackets do the work, there being no comma
to write.} A tuple of two is the \newterm{pair} we want.  Given a pair, a part
comes out by its position---\code{x[0]} is the first, \code{x[1]} the second.
The square brackets say ``the part at this position,'' and the positions are
counted from zero.\footnote{The original's language provides pairs through
three primitive procedures: \code{cons}, which takes two arguments and returns a pair, and
\code{car} and \code{cdr}, which return the first and second parts.  The name
\code{cons} stands for ``construct.''  \code{car} and \code{cdr} derive from
the original implementation of Lisp on the \acronym{IBM 704}, a machine whose
addressing scheme allowed one to reference the ``address'' and ``decrement''
parts of a memory location: \code{car} stands for ``Contents of Address part of
Register'' and \code{cdr} (pronounced ``could-er'') for ``Contents of Decrement
part of Register.''  We shall meet these three again in
\link{Section 2.1.3}, under circumstances that make them worth having.}

\begin{codeBlock}
>>> x = (1, 2)
>>> x[0]
~\textit{1}~
>>> x[1]
~\textit{2}~
\end{codeBlock}

\noindent
Notice that a pair is a data object that can be given a name and manipulated,
just like a primitive data object.  Two further things are worth noticing at
once.  The first is that a pair, once made, cannot be altered: there is no way
to write into \code{x} and change its first part to something else.  This is
the discipline we said at the start of the chapter we would keep, and here the
language keeps it for us.  The second is that a pair may hold pairs, and so on:

\begin{codeBlock}
>>> x = (1, 2)
>>> y = (3, 4)
>>> z = (x, y)
>>> z[0][0]
~\textit{1}~
>>> z[1][0]
~\textit{3}~
\end{codeBlock}

\noindent
In \link{Section 2.2} we will see how this ability to combine pairs means that
pairs can be used as general-purpose building blocks to create all sorts of
complex data structures.  The single compound-data primitive \newterm{pair} is
the only glue we need.  Data objects constructed from pairs are called
\newterm{list-structured} data.

\subsubsection*{Representing rational numbers}

Pairs offer a natural way to complete the rational-number system.  Simply
represent a rational number as a pair of two integers: a numerator and a
denominator.  Then \code{make\_rat}, \code{numer}, and \code{denom} are readily
implemented as follows:\footnote{Another way to define the selectors and
constructor is to name the existing procedures directly, which in the original
reads

\begin{compactCodeBlock}
(define make-rat cons)
(define numer car)
(define denom cdr)
\end{compactCodeBlock}

and is possible because \code{cons}, \code{car} and \code{cdr} are ordinary
procedures, and a procedure is a value that a name may be bound to.  Python
manages half of it.  For the selectors there is the \code{operator} module we
met in \link{Exercise 1.4}, which supplies a procedure for what is ordinarily
written as syntax: as \code{operator.add} is the procedure behind \code{+}, so
\code{operator.itemgetter(0)} is the procedure behind \code{x[0]}, and
\code{numer = operator.itemgetter(0)} is the same move as
\code{(define numer car)}.  The constructor has no such procedure behind it.
Making a pair is written \code{(n, d)}, and there is nothing in the language
that does it---\code{tuple(n, d)} is an error---so one must write
\code{make\_rat = lambda n, d: (n, d)}, which is not naming an existing
procedure but writing a new one.  This is the asymmetry of
\link{Section 1.1.1} again, where \code{+} proved to be syntax rather than a
name; it holds for building data as much as for arithmetic.

The first of those definitions associates the name \code{make-rat} with the
value of the expression \code{cons}, which is the primitive procedure that
constructs pairs; thus \code{make-rat} and \code{cons} become two names for one
procedure.

Defining selectors in this way is efficient, in either language: instead of
\code{numer} \emph{calling} something, \code{numer} \emph{is} that something,
so there is one procedure call where there were two.  On the other hand it
defeats debugging aids that trace calls or set breakpoints on them: you may
want to watch \code{numer} being called, but you certainly do not want to watch
every use of \code{x[0]}.

We have chosen not to use this style of definition in this book.}

\begin{codeBlock}
>>> def make_rat(n, d):
...     return (n, d)

>>> def numer(x):
...     return x[0]

>>> def denom(x):
...     return x[1]
\end{codeBlock}

\noindent
Also, in order to display the results of our computations, we can print
rational numbers by printing the numerator, a slash, and the
denominator:\footnote{\code{print} writes its arguments and then starts a new
line.  The braces in the string mark places where a value is to be put in, and
the \code{f} before the opening quotation mark is what tells Python to do it;
such a string is called an \newterm{f-string}.  \code{print} returns no useful
value, so in the uses of \code{print\_rat} below we show only what it prints,
not what the interpreter reports as the value of the call.}

\begin{codeBlock}
>>> def print_rat(x):
...     print(f"{numer(x)}/{denom(x)}")
\end{codeBlock}

\noindent
Now we can try our rational-number procedures:

\begin{codeBlock}
>>> one_half = make_rat(1, 2)
>>> print_rat(one_half)
~\textit{1/2}~
>>> one_third = make_rat(1, 3)
>>> print_rat(add_rat(one_half, one_third))
~\textit{5/6}~
>>> print_rat(mul_rat(one_half, one_third))
~\textit{1/6}~
>>> print_rat(add_rat(one_third, one_third))
~\textit{6/9}~
\end{codeBlock}

\noindent
As the final example shows, our rational-number implementation does not reduce
rational numbers to lowest terms.  We can remedy this by changing
\code{make\_rat}. If we have a \code{gcd} procedure like the one in
\link{Section 1.2.5} that produces the greatest common divisor of two integers, we can
use \code{gcd} to reduce the numerator and the denominator to lowest terms
before constructing the pair:\footnote{We import one here rather than write it
out.  \code{gcd} in the \code{math} module computes the greatest common divisor
of two integers, which is what its name says and what the procedure of
\link{Section 1.2.5} does.  Which of the two we use makes no difference to
anything in this section: \code{make\_rat} is not affected by where its
\code{gcd} came from, only by what it computes.  That is precisely the black
box of \link{Section 1.1.8}, and it is worth noticing that the barrier works in
this direction too---not only hiding our representation from its users, but
hiding someone else's implementation from us.}

\begin{codeBlock}
>>> from math import gcd

>>> def make_rat(n, d):
...     g = gcd(n, d)
...     return (n // g, d // g)
\end{codeBlock}

\noindent
Now we have

\begin{codeBlock}
>>> print_rat(add_rat(one_third, one_third))
~\textit{2/3}~
\end{codeBlock}

\noindent
as desired.  This modification was accomplished by changing the constructor
\code{make\_rat} without changing any of the procedures (such as \code{add\_rat}
and \code{mul\_rat}) that implement the actual operations.

\begin{exerciseGroup}
\phantomsection\label{Exercise 2.1}
\exerciseHeading{Exercise 2.1:} Define a better version of
\code{make\_rat} that handles both positive and negative arguments.
\code{make\_rat} should normalize the sign so that if the rational number is
positive, both the numerator and denominator are positive, and if the rational
number is negative, only the numerator is negative.
\end{exerciseGroup}

